% Generated draft: requires mathematical and source review.
\documentclass[11pt]{article}
\usepackage[margin=25mm]{geometry}
\usepackage{fontspec}
\setmainfont{Latin Modern Roman}
\usepackage{amsmath,amssymb,mathtools,unicode-math}
\setmathfont{Latin Modern Math}
\usepackage{enumitem,needspace,xurl,hyperref}
\setlength{\parindent}{0pt}
\setlength{\parskip}{5pt}
\setlength{\emergencystretch}{3em}
\newcommand{\N}{\mathbb{N}}
\newcommand{\Z}{\mathbb{Z}}
\newcommand{\Q}{\mathbb{Q}}
\newcommand{\R}{\mathbb{R}}
\newcommand{\C}{\mathbb{C}}
\newcommand{\F}{\mathbb{F}}
\newcommand{\A}{\mathbb{A}}
\newcommand{\PP}{\mathbb{P}}
\newcommand{\E}{\mathbb{E}}
\newcommand{\eps}{\varepsilon}
\newcommand{\dd}{\,\mathrm{d}}
\begin{document}
\section*{UnsolvedMath ID: 701544}
\textbf{Source problem code:} catalogue entry
\subsection{Definitions and mathematical statement}
Fix an algebraic closure $\overline{\Q}$ of $\Q$. For a finite set $S$ of rational primes, let $M_S\subseteq\overline{\Q}$ be the compositum of all finite Galois extensions $L/\Q$ whose degrees are powers of $2$ and which are unramified at every finite prime outside $S$. Thus $M_S$ is the maximal Galois $2$-extension of $\Q$ unramified outside $S\cup\{\infty\}$; ramification at the infinite place is allowed, and the prime $2$ is permitted in $S$.

The source refers to the ``absolute Galois group of the maximal $2$-extension.'' Its literal meaning and the relative Galois group of that extension are, respectively,
\[
 H_S^{\mathrm{abs}}=\operatorname{Gal}(\overline{\Q}/M_S),
 \qquad
 H_S^{\mathrm{rel}}=\operatorname{Gal}(M_S/\Q).
\]
These are different groups. The supplied record does not clarify whether its wording is literal or whether the relative group is intended.

For an explicitly chosen interpretation $H_S$, the precise question is: for every integer $n\geq 0$ and every finite group $G$ of order $2^n$, does there exist a finite set $S$ of rational primes and a continuous surjective homomorphism
\[
 \pi:H_S\twoheadrightarrow G,
 \qquad |S|\leq n,
\]
where $G$ has the discrete topology? The literal formulation takes $H_S=H_S^{\mathrm{abs}}$. The alternative formulation takes $H_S=H_S^{\mathrm{rel}}$; equivalently, it asks for a finite Galois extension $K/\Q$ with $\operatorname{Gal}(K/\Q)\cong G$ that is unramified at every finite prime outside some set $S$ with $|S|\leq n$. The intended choice remains unresolved in this draft.

The supplied source further states that the corresponding assertion for finite $p$-groups, with $p$ odd, is true, as noted by Serre. This is retained as a source-attributed assertion; the record gives neither a separate citation for it nor a clarification of the same Galois-group terminology.
\subsection{Short English statement}
The question asks whether every finite group of order $2^n$ can occur as a continuous finite quotient of a specified Galois group using at most $n$ finite ramification primes, with ramification at infinity allowed. The source's wording leaves unresolved whether the group is the absolute Galois group of the maximal extension or its Galois group over $\Q$. The supplied open label and the odd-prime assertion attributed to Serre are not independently verified here.
\subsection{Sources}
\begin{itemize}
\item Local UnsolvedMath registry, permanent ID $701544$, exact source problem code \texttt{catalogue entry}, title \emph{Kourovka Notebook Problem 18.116}. The supplied registry labels the problem open and records an update date of 2026-10-16; this is supplied registry information, not independent literature verification.
\item \emph{The Kourovka Notebook (Unsolved Problems in Group Theory)}, Problem 18.116, proposed by P. Schmid in 2014. The supplied background identifies the 18th issue (2014) and the 21st issue edition (2026), \texttt{arXiv:1401.0300v48}. The supplied collection description names E. I. Khukhro and V. D. Mazurov as editors. This bibliographic attribution is supplied evidence; the cited text has not been independently consulted.
\item The statement attributes the analogous assertion for odd primes to Serre, without supplying a separate bibliographic reference.
\end{itemize}
\end{document}
